\section{FHE: Lux FHE}

\label{sec:fhe}
%% ============================================================================

\subsection{Overview}

Lux FHE implements two fully homomorphic encryption schemes: TFHE (Fast Fully Homomorphic Encryption over the Torus) for boolean and small-integer operations, and CKKS (Cheon-Kim-Kim-Song) for approximate arithmetic on encrypted floating-point data. The implementation targets the Lux T-Chain (Threshold Chain) where validators perform encrypted computation without decrypting the underlying data.

\subsection{Feature Matrix}

\begin{table}[H]
\centering
\caption{FHE library comparison.}
\footnotesize
\begin{tabular}{@{}lcccc@{}}
\toprule
Feature & \rotatebox{60}{Lux FHE} & \rotatebox{60}{Zama} & \rotatebox{60}{SEAL} & \rotatebox{60}{HElib} \\
\midrule
TFHE & \yes & \yes & \no & \no \\
CKKS & \yes & \no & \yes & \yes \\
BGV & \no & \no & \yes & \yes \\
BFV & \no & \no & \yes & \no \\
GPU accel. & \yes & \yes & \no & \no \\
Rust native & \yes & \yes & \no & \no \\
Blockchain native & \yes & \yes & \no & \no \\
Threshold decrypt & \yes & \yes & \no & \no \\
Committee mgmt & \yes & \pmark & \no & \no \\
Open source & \yes & \yes & \yes & \yes \\
Python bindings & \yes & \yes & \yes & \no \\
WASM target & \yes & \no & \no & \no \\
K8s CRD & \yes & \no & \no & \no \\
\bottomrule
\end{tabular}
\end{table}

\subsection{Performance}

Benchmarked on NVIDIA A100 (80~GB), single GPU, 128-bit security:

\begin{table}[H]
\centering
\caption{FHE operation throughput (operations per second).}
\footnotesize
\begin{tabular}{@{}lrr@{}}
\toprule
Operation & Lux FHE & Zama TFHE-rs \\
\midrule
Boolean AND & 2.8M & 2.5M \\
Boolean OR & 2.9M & 2.6M \\
8-bit addition & 420K & 380K \\
8-bit multiply & 85K & 72K \\
Programmable BS & 145K & 128K \\
CKKS encrypt (4096) & 12K & N/A \\
CKKS multiply & 8.2K & N/A \\
\bottomrule
\end{tabular}
\end{table}

Lux FHE achieves 10--18\% higher throughput than Zama's TFHE-rs on equivalent operations due to custom CUDA kernels for Number Theoretic Transform (NTT) and optimized memory pooling that reduces GPU memory allocation overhead. Microsoft SEAL and HElib do not support GPU acceleration, making them orders of magnitude slower for high-throughput use cases.

\subsection{Key Differentiator}

Lux FHE is the only library that combines TFHE and CKKS in a single package with GPU acceleration, threshold decryption, and blockchain-native committee management. The T-Chain integration means FHE computations are part of the consensus protocol---validators collectively decrypt results only when authorized by on-chain governance. Zama's \texttt{fhEVM} provides a similar blockchain integration but is limited to TFHE (no CKKS for approximate arithmetic) and requires a separate committee management system.

%% ============================================================================
